\chapter{Auswertung, Fehlerrechnung und Diskussion der Messergebnisse}
\section{Energiekalibration}
In der ersten Aufgabe wurde eine Energiekalibrierung des Detektors durchgeführt, um einen Zusammenhang zwischen Kanalnummern und den tatsächlichen Energien der detektierten Photonen herzustellen. Hierzu wurden bekannte $\gamma$-Peaks der Referenzquellen $^{22}$Na, $^{57}$Co und $^{137}$Cs verwendet. Die Positionen dieser Peaks im Spektrum wurden jeweils durch Fitten mit einer Gaußfunktion bestimmt.

\begin{table}[H]
\centering
\begin{tabular}{c|c|c|c}
Probe & Energie in keV & Kanal (NaI) & Kanal (Ge) \\
\hline
$^{22}$Na     & 511      & ...  & ... \\
$^{22}$Na     & 1275     & ...  & ... \\
$^{137}$Cs    & 662      & ...  & ... \\
$^{57}$Co     & 122      & ...  & ... \\
\end{tabular}
\caption{Literaturwerte der $\gamma$-Energien und gemessene Kanalnummern der bekannten Peaks.}
\end{table}


Anschließend wurden die gemessenen Kanalpositionen der Peaks den entsprechenden Literaturwerten der $\gamma$-Energien zugeordnet. Aus diesen Datenpunkten wurde ein linearer Fit durchgeführt, der die Kalibrierfunktion beschreibt. Diese liefert die Parameter des linearen Zusammenhangs zwischen Kanalnummer $K$ und Energie $E$:

\[
E_{NaJ} = (0,0\pm 0,0) \,\frac{\text{keV}}{\text{Kanal}}\cdot K + (0,0\pm 0,0) \,\text{keV}
\]

\[
E_{Ge} = (0,0\pm 0,0) \,\frac{\text{keV}}{\text{Kanal}}\cdot K + (0,0\pm 0,0) \,\text{keV}
\]

Die so erhaltene Kalibrierfunktion ermöglicht es, künftig aus der gemessenen Kanalnummer eines Peaks auf die zugehörige Photonenenergie zu schließen.

\section{Co60-Spektrum}

\begin{table}[H]
\centering
\begin{tabular}{c|c|c}
Peak & Energie in keV (NaJ) & Energie in keV (Ge) \\
\hline
Rückstreupeak 1   & ...      & ...   \\
Rückstreupeak 2    & ...     & ...   \\
Comptonkante 1    & ...    & ...  \\
Comptonkante 2    & ...      & ...   \\
Photopeak 1   &...      & ...   \\
Photopeak 2   & ...     & ...   \\
\end{tabular}
\caption{Energien für Rückstreu -und Photopeaks, sowie Comptonkante für beide Detektoren}
\end{table}




\section{Koinzidenzanalyse}